\section{Introduction}\label{sec:intro}
A market simulator is useful only through what its generated paths imply. The distribution of returns matters, but so do the probability of receiving a fill, the time needed to exhaust a queue, the price paid for a specified quantity, and the response to an action selected by a trading policy. Small errors in a displayed book representation need not imply small errors in these quantities. Depleting the final lot can change the best quote. A reserve replenishment can prevent that change. Two states with the same visible depth can therefore have different continuation laws.

Generative neural operators offer a natural representation for financial objects indexed by price, maturity, asset, or time. Their function-space formulation separates the probabilistic object from the grid on which it is computed \citep{kovachki2023,rahman2022,kerrigan2024functional}. The representation alone does not ensure that generated events are valid or that a simulated decision has approximately the right value. Those properties require a connection between field geometry, latent information, event dynamics, and decision loss.

We develop that connection in a partially observed market model. The latent state has a conditional law $\pi_t$ on a separable Hilbert space. It can describe liquidity-regime factors, unresolved order-flow intensities, and transformed reserve coordinates. It is not automatically an identifiable physical inventory. A conditional generative operator approximates $\pi_t$, and a second operator maps each latent state to a field of event-intensity densities. The densities are integrated over physical event cells and supplied to an admissible event simulator.

The central bound has the form
\begin{equation}\label{eq:intro-bound}
 \TV(\mathbb P^a_{[0,T]},\widehat{\mathbb P}^a_{[0,T]})
 \le \min\left\{1,\,
 \sqrt M\,\E^a\int_0^T
 \left[r_t+K_t\sqrt{\tau_{t,m}^2+d_{t,m}^2}\right]\dd t\right\}.
\end{equation}
Here $a$ is a common history-dependent policy, $M$ is the mass of the reference measure on event space, $r_t$ is the target-averaged conditional intensity-field error, $K_t$ is a learned latent-to-intensity sensitivity, $\tau_{t,m}$ is the unresolved posterior tail, and $d_{t,m}$ is a conditional refinement certificate. All quantities are evaluated along true histories under that policy. In particular, the theorem does not replace those histories by a passive training distribution without justification.

The error criterion in \eqref{eq:intro-bound} is strong enough to handle discontinuous path functionals. An event probability changes by at most the right-hand side. A cost taking values in $[0,B]$ changes in expectation by at most $B$ times that bound. Uniformity over a policy class gives a $2B$ multiple for model-selection regret. Tail risk introduces the explicit amplification factor $(1-\alpha)^{-1}$ for $\CVaR_\alpha$. The same guarantee does not imply a continuous value-at-risk estimate at atoms.

For decisions with a specified objective, a full-path bound can be unnecessarily conservative. We derive a continuation-value sensitivity identity and connect it to rate confidence sequences constructed from observed counts and exposures. In a finite observable state model these intervals remain valid under adaptive actions and stopping; robust dynamic programming converts the same event into a guarantee for the selected execution policy. This specialization requires observed state-action exposure and does not infer hidden reserve exposure from displayed data. Its purpose is to supply a complete statistical-to-decision chain where identification is explicit.

\subsection{Relation to existing work}
Queue-based Markov models supply a tractable mathematical description of order-book events \citep{cont2010}. Learned world-agent simulators generate market activity from data \citep{coletta2022}, while LOB-Bench evaluates distributions, conditional statistics, and event-response behavior in generated message streams \citep{nagy2025}. Our aim is to identify sufficient conditions under which such distributional approximation controls financial decisions. The paper does not claim to introduce event-driven simulation, generative market models, or hidden-state filtering.

The underlying methods are also established. Conditional transport and function-space flow matching provide generative representations \citep{hosseini2025,kerrigan2024functional}; disintegration and conditional probability supply latent posteriors \citep{kallenberg2021}; and common-Poisson coupling is a standard point-process comparison method \citep{bremaud1981}. Our contribution is the combined certificate through cell-integrated intensity fields, with an explicit separation between posterior resolution, conditional fitting, and sensitivity to latent errors. This separation makes the numerical-resolution constant and the policy-distribution requirement visible.

Financial admissibility has more than one meaning. Valid matching events concern the support of a physical market simulator. Martingale pricing concerns a separately specified pricing measure. Neural-SDE market models impose financial constraints on option dynamics \citep{cohen2021}, and VolGAN combines generative volatility modeling with scenario treatment for arbitrage constraints \citep{vuletic2024}. Risk-neutral neural operators for joint SPX--VIX term structures have also been proposed \citep{zhang2025}. We provide a simpler positive martingale construction to show exactly where an operator can enter a pricing model and to derive an adapted volatility-to-price bound. We do not identify the pricing law with the historical law or claim that a generated volatility field calibrates observed option prices.

\subsection{Main results and scope}
The principal results are: a structural validity proposition for event generation; a posterior-intensity identity and a quantitative obstruction to snapshot-only prediction; a field-to-path comparison theorem; a multiresolution posterior bound; execution and risk guarantees; a liquidity-quantile transport inequality; and an adapted martingale construction. The decision-specific extension adds event-rate confidence sequences, a robust policy sandwich, and a numerical residual allowance. A bicausal stopping comparison treats early exercise and distinguishes information timing from terminal calibration. Each result states the assumptions that make its claim meaningful.

The numerical studies include a complete reserve model with a trained rate network, exact survival-aware filtering, two exact thinning implementations, and an optimized finite class of execution deadlines. Matrix exponentiation evaluates policy errors independently of simulation. A separate queue experiment compares coupling and likelihood certificates, including an informative 99-percent tail-risk regime. A three-stream observed-event study fits rates without oracle rate labels, computes confidence boxes, and compares nominal and robust policies across three data budgets. These are controlled synthetic validations rather than exchange-data calibration or measurements of undisplayed exchange inventory. Formal verification is similarly scoped: the Lean files prove their stated mathematical declarations, and a coverage map identifies which analytic results are proved in this manuscript but not mechanized.
